\pdfoutput=1
\pdfinfoomitdate=1
\pdftrailerid{}
\pdfsuppressptexinfo=15
\documentclass[indonesian]{shinybook}
\usepackage{tikz}
\usepackage{pgfplots}
\usepgfplotslibrary{groupplots}
\pgfplotsset{compat=1.18}
\usetikzlibrary{intersections,calc}
\usetikzlibrary{arrows.meta, decorations.markings}
\usepackage{glossaries-extra}
\usepackage{needspace}

\newcommand{\half}{\frac{1}{2}}
\newcommand{\N}{\mathbb{N}}
\newcommand{\Z}{\mathbb{Z}}
\newcommand{\Q}{\mathbb{Q}}
\newcommand{\R}{\mathbb{R}}
\newcommand{\Rb}{\overline{\mathbb{R}}}
\newcommand{\bR}{\overline{\mathbb{R}}}
\newcommand{\F}{\mathbb{F}}
\newcommand{\1}{\mathbb{1}}
\newcommand{\linop}{\mathcal{L}}
\newcommand{\Rnn}{\mathbb{R}^{n\times n}}
\newcommand{\Cnn}{\mathbb{C}^{n\times n}}
\newcommand{\E}{\mathbb{E}}
\usepackage{pgffor}
\foreach \x in {A,...,Z}{%
    \expandafter\xdef\csname cal\x\endcsname{\noexpand\ensuremath{\noexpand\mathcal{\x}}}
}

\renewcommand{\implies}{\Rightarrow}
\newcommand{\setto}{\rightrightarrows}
\newcommand{\equivalent}{\Leftrightarrow}
\newcommand{\eps}{\varepsilon}
\renewcommand{\phi}{\varphi}
\newcommand{\supp}{\operatorname{\mathrm{supp}}}
\renewcommand{\div}{\operatorname{\mathrm{div}}}
\newcommand{\dom}{\operatorname{\mathrm{dom}}}
\renewcommand{\ker}{\operatorname{\mathrm{ker}}}
\newcommand{\conv}{\operatorname{\mathrm{conv}}}
\newcommand{\rg}{\operatorname{\mathrm{rg}}}
\newcommand{\graph}{\operatorname{\mathrm{graph}}}
\newcommand{\tr}{\operatorname{\mathrm{tr}}}
\newcommand{\epi}{\operatorname{\mathrm{epi}}}
% Never redefine TeX's primitive \span: amsmath alignment environments use it.
% Use \spann for the linear-span operator in reader prose and formulae.
\newcommand{\sign}{\operatorname{\mathrm{sign}}}
\DeclareMathOperator{\Id}{\mathrm{Id}}
\newcommand{\prox}{\mathrm{prox}}
\newcommand{\proj}{\mathrm{proj}}
\newcommand{\spann}{\operatorname{\mathrm{span}}}
\newcommand{\Lloc}{L^1_{\mathrm{loc}}}
\newcommand{\Cinf}{C^\infty_0}

% Nama lingkungan dilokalkan; penomoran dan topologi sumber dipertahankan.
\newtheorem{theorem}{Teorema}[chapter]
\newtheorem{lemma}[theorem]{Lema}
\newtheorem{cor}[theorem]{Korolari}
\theoremstyle{definition}
\newtheorem{defn}[theorem]{Definisi}
\newtheorem{prop}[theorem]{Proposisi}
\newtheorem{example}[theorem]{Contoh}
\newtheorem{assumption}[theorem]{Asumsi}
\newtheorem{exercise}[theorem]{Latihan}
\newtheorem*{rem}{Catatan}

\crefname{theorem}{teorema}{teorema}
\Crefname{theorem}{Teorema}{Teorema}
\crefname{lemma}{lema}{lema}
\Crefname{lemma}{Lema}{Lema}
\crefname{cor}{korolari}{korolari}
\Crefname{cor}{Korolari}{Korolari}
\crefname{prop}{proposisi}{proposisi}
\Crefname{prop}{Proposisi}{Proposisi}
\crefname{defn}{definisi}{definisi}
\Crefname{defn}{Definisi}{Definisi}
\crefname{example}{contoh}{contoh}
\Crefname{example}{Contoh}{Contoh}
\crefname{assumption}{asumsi}{asumsi}
\Crefname{assumption}{Asumsi}{Asumsi}
\crefname{rem}{catatan}{catatan}
\Crefname{rem}{Catatan}{Catatan}
\crefname{exercise}{latihan}{latihan}
\Crefname{exercise}{Latihan}{Latihan}

\renewcommand{\proofname}{Bukti}
\renewcommand{\contentsname}{Daftar Isi}
\renewcommand{\figurename}{Gambar}

\newcommand{\norm}[1]{\|#1\|}
\newcommand{\abs}[1]{\left|#1\right|}
\newcommand{\inner}[2]{\left\langle#1 , #2\right\rangle}
\newcommand{\dual}[1]{\langle #1 \rangle}
\newcommand{\setof}[2]{\left\{#1 \;\middle|\; #2\right\}}
\newcommand{\wkto}{\rightharpoonup}
\DeclareMathOperator{\Exp}{\mathbb{E}}
\DeclareMathOperator{\Var}{\mathbb{V}}
\DeclareMathOperator{\Cov}{\mathrm{Cov}}
\newcommand{\Prob}{\mathbb{P}}

\usepackage{enumerate}
\newcommand{\weaki}{\protect{\u{\i}}}
\AtBeginEnvironment{remark}{\small}

\ifdefined\OIntegratedReader
  % Keep examples as ordinary theorem blocks in the tagged master. The legacy
  % mdframed/tablefootnote patches interfere with the synchronized latex-lab
  % hooks, while no canonical body calls \tablefootnote.
\else
  \usepackage{mdframed}
  \mdfdefinestyle{examplebg}{
      backgroundcolor=structure!7,%
      hidealllines=true,%
      splittopskip=\topskip,
      frametitleaboveskip=0,
      frametitlebelowskip=0,
      innertopmargin=-.3em,%
      innerbottommargin=.7em,%
      innerleftmargin=.7em,%
      innerrightmargin=.7em,%
  }
  \surroundwithmdframed[style=examplebg]{example}
  \usepackage{tablefootnote}
  \makeatletter
  \AfterEndEnvironment{mdframed}{%
   \tfn@tablefootnoteprintout%
   \gdef\tfn@fnt{0}%
  }
  \makeatother
\fi

\usepackage{enumitem}
\setlist[enumerate]{label=(\roman*)}
\usepackage{scrhack}

\newcommand{\ie}{\textit{yaitu}}
\newcommand{\eg}{\textit{misalnya}}
\newcommand{\cf}{\textit{bandingkan}}
\newcommand{\emptyarg}{\,\cdot\,}
\newcommand{\dd}{\mathrm{d}}
\newcommand{\todo}[1]{{\color{red}Perlu dikerjakan: #1}}
\newcommand{\Oc}{\mathcal{O}}
\newcommand{\Lc}{\mathcal{L}}
\newcommand{\Gc}{\mathcal{G}}
\newcommand{\Ic}{\mathcal{I}}
\newcommand{\Pc}{\mathcal{P}}
\newcommand{\Mc}{\mathcal{M}}
\renewcommand{\epsilon}{\varepsilon}
\newcommand{\diag}{\mathrm{diag}}

\addbibresource{../../authority/habring/source-v1/references.bib}
\microtypesetup{expansion=false}

% Lokalisasi ini sengaja hanya berlaku pada unit mandiri Bab 7.
\crefname{equation}{persamaan}{persamaan}
\Crefname{equation}{Persamaan}{Persamaan}

\title{Optimisasi Konveks}
\subtitle{Unit 7: Dualitas\\Edisi Bahasa Indonesia}
\author{Andreas Habring}
\date{Sumber edisi v1: 13 Juli 2026}
\publishers{\small
Terjemahan kerja bahasa Indonesia dari \emph{Lecture Notes: Convex Optimization}, arXiv:2607.11664v1.\\
Sumber asli dan terjemahan ini dilisensikan CC BY 4.0. Perubahan matematis dijelaskan dalam catatan koreksi.\\
Terjemahan mandiri ini bukan karya resmi atau dukungan Andreas Habring maupun TU Graz.}
\hypersetup{
    pdftitle={Optimisasi Konveks — Unit 7: Dualitas — Edisi Bahasa Indonesia},
    pdfauthor={Andreas Habring},
    pdflang={id-ID},
}

\begin{document}
\maketitle

\frontmatter
\chapter*{Tentang unit ini}
Unit ini menerjemahkan secara utuh Bab 7, ``Duality'', dari Andreas Habring, \emph{Lecture Notes: Convex Optimization}, arXiv:2607.11664v1. Sumber berwenang untuk unit ini adalah berkas \texttt{duality.tex} edisi v1 dengan SHA-256 berikut:
\begin{center}
\small\texttt{0b112dee2582813cec5629c02df1dda32}\\[-0.2em]
\texttt{9f690f944b60f4694b1c5762129bea9}
\end{center}

Bab-bab prasyarat tidak diulang. Pembaca diharapkan telah mengenal ruang bernorma dan ruang Hilbert, dual dan adjoin operator, semikontinuitas bawah, epigraf dan pemisahan Hahn--Banach, subgradien, operator proksimal, serta dasar analisis konvergensi. Seluruh definisi, contoh, lema, teorema, bukti, latihan, sitasi, dan catatan kaki dalam bab sumber dipertahankan dalam urutan sumber. Kekeliruan yang dapat ditentukan secara matematis diperbaiki secara terbuka dan diringkas pada akhir unit.

Edisi ini berdiri sendiri untuk dibaca dan dibangun, tetapi tidak mengklaim menggantikan konteks lengkap catatan kuliah asal. Istilah teknis Indonesia dipakai secara konsisten; singkatan PDHG dan ADMM dipertahankan karena merupakan nama baku algoritme.

\tableofcontents
\mainmatter
\setcounter{chapter}{6}
% H07-S001 | sumber authority/habring/source-v1/duality.tex baris 1--24
% segment-id: d90.hab.v1.ch07.seg0001
\chapter{Dualitas}\label{chapter:duality}
Dalam bab ini, semua ruang vektor dipahami sebagai ruang Hilbert riil berdimensi hingga. Dual yang digunakan adalah dual kontinu; melalui representasi Riesz kita mengidentifikasi $V^*\cong V$ dan, khususnya, $V^{**}\cong V$. Untuk setiap ruang demikian $H$, tuliskan $\Gamma_0(H)$ bagi kelas fungsi $H\rightarrow(-\infty,+\infty]$ yang proper, konveks, dan semikontinu bawah.

\section{Dualitas Fenchel}

\begin{defn}[Konjugat Fenchel/konveks]
    Misalkan $f:V\rightarrow(-\infty,+\infty]$ proper. Konjugat Fenchel, atau konjugat konveks, didefinisikan sebagai
    \begin{equation}
        \begin{aligned}
            f^*:V^*&\rightarrow(-\infty,+\infty]\\
            f^*(x^*) &= \sup_{x\in V} \inner{x^*}{x} - f(x).
        \end{aligned}
    \end{equation}
\end{defn}

\begin{example}
    \begin{enumerate}
        \item Fungsi indikator: Misalkan $f=\delta_C$ dengan $C\subset V$ tak kosong dan konveks. Maka $f^*(x^*) = \sup_{x\in C}\inner{x^*}{x}$ adalah apa yang disebut fungsi pendukung.
        \item Norma: Misalkan $f = \|\emptyarg\|$ adalah suatu norma. Maka $f^* = \delta_{\overline{B}_{\|\emptyarg\|_*,1}(0)}$. Artinya, $f^*$ adalah fungsi indikator pada bola satuan tertutup terhadap norma dual.
        \item Sebaliknya, konjugat dari $f = \delta_{\overline{B}_{\|\emptyarg\|_*,1}(0)}$ adalah $f^* = \|\emptyarg\|$ setelah identifikasi $V^{**}\cong V$.
    \end{enumerate}
\end{example}

% \todo{Alex: Kalkulus sebagai latihan}
% H07-S002 | sumber authority/habring/source-v1/duality.tex baris 25--60
% segment-id: d90.hab.v1.ch07.seg0002
\begin{lemma}
    Misalkan $f$ proper dan konveks, serta andaikan bahwa $\partial f(\hat x) \neq \emptyset$ untuk suatu $\hat x\in V$. Dalam dimensi hingga, syarat tambahan ini otomatis dipenuhi pada suatu titik di interior relatif domain efektif; \cf~\cite[Corollary 3.19]{beck2017first}. Maka $f^*$ proper, konveks, dan semikontinu bawah (lsc).
\end{lemma}
\begin{proof}
    Karena $f$ proper, terdapat $x_0\in V$ sedemikian sehingga $f(x_0)\in \R$. Akibatnya, untuk setiap $x^*$,
    \begin{equation}
        f^*(x^*)\geq \inner{x^*}{x_0} - f(x_0)> -\infty.
    \end{equation}
    Ambil $g\in \partial f(\hat x)$. Berdasarkan definisi subgradien, kita mempunyai
    \begin{equation}
        f(\hat x) + \inner{g}{x-\hat x}\leq f(x)
    \end{equation}
    dan, dengan demikian,
    \begin{equation}
        \begin{aligned}
            f^*(g)
            =& \sup_{x\in V}\inner{g}{x} - f(x)\\
            \leq& \sup_{x\in V}f(x) - f(\hat x) + \inner{g}{\hat x} - f(x)\\
            \leq& \sup_{x\in V}- f(\hat x) + \inner{g}{\hat x} <\infty.
        \end{aligned}
    \end{equation}
    Kekonveksan dan semikontinuitas bawah langsung mengikuti fakta bahwa konjugat konveks merupakan supremum titik demi titik dari fungsi-fungsi afin kontinu, yang karena itu konveks dan lsc\footnote{Di sini kekonveksan telah ditunjukkan secara langsung. Rincian untuk lsc ditinggalkan sebagai latihan.}.
\end{proof}

\begin{lemma}[Ketaksamaan Fenchel]
    Misalkan $f$ proper. Maka
    \begin{equation}
        f^*(x^*) + f(x)\geq \inner{x^*}{x}.
    \end{equation}
\end{lemma}
\begin{proof}
    Latihan.
\end{proof}

Untuk fungsi proper $f$ dengan $f^*$ proper, proses pengambilan konjugat dapat kita \emph{iterasikan} untuk memperoleh \emph{bikonjugat}. Kondisi ini dijamin oleh lema sebelumnya pada kasus konveks yang memenuhi hipotesisnya. Demi kesederhanaan, kita tulis $f^{**} \coloneq (f^*)^*$. Karena diasumsikan $V^{**}=V$, bikonjugat kembali terdefinisi pada ruang asal $V$.

% H07-S003 | sumber authority/habring/source-v1/duality.tex baris 61--141
% segment-id: d90.hab.v1.ch07.seg0003
\begin{lemma}
    Misalkan $f$ dan $f^*$ proper. Untuk setiap $x\in V$ berlaku $f^{**}(x)\leq f(x)$.
\end{lemma}
\begin{proof}
    Pertama, perhatikan bahwa untuk setiap $x$, konjugat memenuhi
    \begin{equation}
        f^*(x^*)\geq \inner{x^*}{x}-f(x).
    \end{equation}
    Oleh karena itu,
    \begin{equation}
        \begin{aligned}
            f^{**}(x)
            =& \sup_{x^*\in V^*}\inner{x^*}{x} - f^*(x^*)\\
            =& \sup_{x^*\in V^*}\inf_y \inner{x^*}{x-y} + f(y)\\
            \leq& \sup_{x^*\in V^*} \inner{x^*}{x-x} + f(x)\\
            =& f(x).
        \end{aligned}
    \end{equation}
\end{proof}

\begin{lemma}
    Misalkan $f$ proper, konveks, dan lsc. Maka $f^{**} = f$.
\end{lemma}
\begin{proof}
    Karena selalu berlaku $f^{**}\leq f$, tinggal ditunjukkan bahwa $f^{**}\geq f$. Andaikan sebaliknya: terdapat $x_0\in V$ sedemikian sehingga
    \begin{equation}
        f^{**}(x_0) < f(x_0).
    \end{equation}
    Ini berarti $(x_0, f^{**}(x_0))\not\in \epi(f)$. Karena $\epi(f)$ tertutup dan konveks berdasarkan asumsi, teorema Hahn--Banach dapat diterapkan untuk mendapatkan $(g,\alpha)\in V^*\times \R$ dan $c_1,c_2\in \R$ sedemikian sehingga
    \begin{equation}
        \inner{(g,\alpha)}{(x_0,f^{**}(x_0))}<c_1<c_2<\inner{(g,\alpha)}{(x,t)},\quad (x,t)\in \epi(f).
    \end{equation}
    Karena $f$ proper, pilih $\bar x\in\dom(f)$. Dengan memasukkan $(\bar x,t)\in\epi(f)$ dan membiarkan $t\rightarrow\infty$, diperoleh $\alpha\geq 0$.
    Sekarang andaikan $\alpha>0$. Kita dapat membagi dengan $\alpha$ dan memperoleh
    \begin{equation}
        \inner{\frac{g}{\alpha}}{x_0-x} + f^{**}(x_0)-f(x)\leq \frac{c_1-c_2}{\alpha}<0.
    \end{equation}
    Ketaksamaan di atas berlaku untuk setiap $x\in \dom(f)$ dan juga secara otomatis untuk setiap $x$ dengan $f(x)=\infty$. Mengambil supremum terhadap $x$ menghasilkan
    \begin{equation}
        \inner{\frac{g}{\alpha}}{x_0} + f^{**}(x_0)+f^*(-\frac{g}{\alpha})<0,
    \end{equation}
    yang bertentangan dengan ketaksamaan Fenchel.
    Pada kasus yang tersisa, yaitu $\alpha=0$, kita mempunyai
    \begin{equation}
        \inner{g}{x_0-x} < c_1-c_2<0.
    \end{equation}
    Ambil sembarang $\hat y\in \dom(f^*)$. Untuk $\hat g = g-\epsilon \hat y$, kita memperoleh
    \begin{equation}
        \begin{aligned}
            \inner{\hat g}{x_0-x} + \epsilon(f^{**}(x_0)-f(x))
            \leq& \inner{g}{x_0-x} + \epsilon( \inner{-\hat y}{x_0-x} + f^{**}(x_0)-f(x))\\
            \leq& c_1-c_2 + \epsilon( \inner{\hat y}{-x_0} + f^{**}(x_0)+f^*(\hat y)).
        \end{aligned}
    \end{equation}
    Perhatikan bahwa $\hat y$ dan $x_0$ tetap. Karena itu, $\epsilon>0$ dapat dipilih cukup kecil agar ruas kanan tetap negatif ketat. Tuliskan $c \coloneq c_1-c_2 + \epsilon( \inner{-\hat y}{x_0} + f^{**}(x_0)+f^*(\hat y))<0$. Setelah membagi dengan $\epsilon$, diperoleh
    \begin{equation}
        \begin{aligned}
            \inner{\frac{\hat g}{\epsilon}}{x_0-x} + f^{**}(x_0)-f(x)
            \leq \frac{c}{\epsilon}<0.
        \end{aligned}
    \end{equation}
    Sekali lagi mengambil supremum terhadap semua $x$ menghasilkan
    \begin{equation}
        \begin{aligned}
            \inner{\frac{\hat g}{\epsilon}}{x_0} + f^{**}(x_0)+f^*(-\frac{\hat g}{\epsilon})
            \leq \frac{c}{\epsilon}<0,
        \end{aligned}
    \end{equation}
    Hal ini bertentangan dengan ketaksamaan Fenchel. Jadi, asumsi $f^{**}(x_0)<f(x_0)$ selalu menimbulkan kontradiksi. Karena itu, $f^{**}(x_0)\geq f(x_0)$, dan bukti selesai.
\end{proof}

\begin{rem}
    Jika $f$ tidak sekaligus konveks dan tertutup, maka, selama keluarga minoran proper di bawah tidak kosong, $f^{**}$ merupakan minoran konveks semikontinu bawah terbesar dari $f$, yang disebut regularisasi $\Gamma$. Lebih tepatnya, tuliskan
    \begin{equation}
        \Gamma(V)=\{\ell:V\rightarrow(-\infty,+\infty]\;|\; \ell\text{ proper, tertutup, dan konveks}\}.
    \end{equation}
    Jika $\{\ell\in\Gamma(V):\ell\leq f\}\neq\emptyset$, maka
    \begin{equation}
        f^{**}(x) = \sup_{\substack{\ell\in \Gamma(V),\\
        \ell\leq f}}\ell (x).
    \end{equation}
    Jika keluarga minoran tersebut kosong, supremum pada ruas kanan dipahami sebagai fungsi yang identik $-\infty$. Fungsi ini tidak termasuk dalam $\Gamma(V)$ dan, di bawah konvensi bab ini yang mendefinisikan konjugasi hanya bagi masukan proper, tidak dinotasikan sebagai $f^{**}$.
\end{rem}
% H07-S004 | sumber authority/habring/source-v1/duality.tex baris 142--186
% segment-id: d90.hab.v1.ch07.seg0004
Konjugat mempunyai hubungan yang erat dengan subgradien, sebagaimana ditunjukkan oleh lema berikut.
\begin{lemma}\label{duality:lemma:subdiff_conjugate}
    Misalkan $f$ proper, konveks, dan tertutup. Pernyataan-pernyataan berikut ekuivalen:
    \begin{enumerate}
        \item $f(x) + f^*(x^*) = \inner{x^*}{x}$.\label{duality:item1}
        \item $x^*\in \partial f(x)$.\label{duality:item2}
        \item $x\in \partial f^*(x^*)$.\label{duality:item3}
    \end{enumerate}
\end{lemma}
\begin{proof}
    \leavevmode\par\smallskip
    \begin{enumerate}
        \item[\ref{duality:item1}$\rightarrow$\ref{duality:item2}:] Pernyataan~\ref{duality:item1} menyiratkan bahwa untuk setiap $z\in V$,
        \begin{equation}
            \begin{aligned}
                -f(x) + \inner{x^*}{x} = f^*(x^*) \geq \inner{x^*}{z} - f(z).
            \end{aligned}
        \end{equation}
        Berdasarkan definisi, ini menyiratkan $x^*\in \partial f(x)$.
        \item[\ref{duality:item2}$\rightarrow$\ref{duality:item1}:]
        Sebaliknya, pernyataan~\ref{duality:item2} menghasilkan, untuk setiap $z$,
        \begin{equation}
            \begin{aligned}
                -f(x) + \inner{x^*}{x} \geq \inner{x^*}{z} - f(z).
            \end{aligned}
        \end{equation}
        Dengan mengambil supremum terhadap $z$, diperoleh $\inner{x^*}{x} \geq f(x) + f^*(x^*)$. Bersama ketaksamaan Fenchel, ini membuktikan hasilnya.
        \item[\ref{duality:item1}$\leftrightarrow$\ref{duality:item3}:]
        Karena $f$ proper, konveks, dan tertutup, pernyataan~\ref{duality:item1} ekuivalen dengan
        \begin{equation}
            f^{**}(x) + f^*(x^*) = \inner{x^*}{x}.
        \end{equation}
        Bukti kemudian diperoleh dengan mengulangi argumen di atas.
    \end{enumerate}
\end{proof}

\begin{theorem}[Identitas Moreau]\label{duality:thm:Moreau:id}
    Misalkan $f\in\Gamma_0(V)$. Maka
    \begin{equation}
        x = \prox_f(x) + \prox_{f^*}(x).
    \end{equation}
\end{theorem}
\begin{proof}
    Pertama, perhatikan bahwa kedua pemetaan proksimal terdefinisi dengan baik dan bernilai tunggal (mengapa?). Demi kesederhanaan, tuliskan $p = \prox_f(x)$. Dari syarat optimalitas diperoleh $x-p\in\partial f(p)$ dan, berdasarkan~\cref{duality:lemma:subdiff_conjugate}, $p = x - (x-p)\in \partial f^*(x-p)$. Pernyataan terakhir menyiratkan $\prox_{f^*}(x) = x-p$, sehingga bukti selesai.
\end{proof}
% H07-S005 | sumber authority/habring/source-v1/duality.tex baris 187--247
% segment-id: d90.hab.v1.ch07.seg0005
\begin{theorem}[Fenchel--Rockafellar]\label{duality:eq:fenchel_duality}
    Misalkan $X$ dan $Y$ ruang Hilbert riil berdimensi hingga, $f\in\Gamma_0(X)$, $g\in\Gamma_0(Y)$, dan $K:X\rightarrow Y$ operator linear terbatas. Selain itu, andaikan terdapat $x_0 \in \dom(f) \cap \dom(g\circ K)$ dengan $K x_0 \in \dom(g)^\circ$. Andaikan masalah (primal)
    \begin{equation}
        \min_{x} f(x) + g(Kx)
    \end{equation}
    mempunyai solusi $\hat x$. Maka masalah dual
    \begin{equation}
        \begin{aligned}
            \max_{x^*} -f^*(-K^*x^*) - g^*(x^*)
        \end{aligned}
    \end{equation}
    juga mempunyai solusi $\hat x^*$, dan berlaku dualitas kuat, yaitu
    \begin{equation}
        \begin{aligned}
            \min_{x} f(x) + g(Kx)
            = \max_{x^*} -f^*(-K^*x^*) - g^*(x^*).
        \end{aligned}
    \end{equation}
    Selain itu, solusi primal dan dual memenuhi
    \begin{equation}
        \begin{cases}
            -K^*\hat x^*&\in \partial f(\hat x),\\
            \hat x^* &\in \partial g(K\hat x).
        \end{cases}
    \end{equation}
\end{theorem}
\begin{proof}
    Pertama, secara langsung diperoleh
    \begin{equation}\label{duality:eq:fenchel_rocka}
        \begin{aligned}
            \inf_{x} f(x) + g(Kx)
            =& \inf_{x} f(x) + g^{**}(Kx)\\
            =& \inf_{x}\sup_{x^*} f(x) + \inner{Kx}{x^*} - g^*(x^*)\\
            \geq & \sup_{x^*} \inf_{x} f(x) + \inner{Kx}{x^*} - g^*(x^*)\\
            = & \sup_{x^*} -\sup_{x} -f(x) + \inner{x}{-K^*x^*} - g^*(x^*)\\
            = & \sup_{x^*} -f^*(-K^*x^*) - g^*(x^*).
        \end{aligned}
    \end{equation}
    Jadi, tinggal dibuktikan ketaksamaan sebaliknya.
    Hipotesis interior di atas memenuhi syarat kualifikasi bagi aturan jumlah dan rantai subdiferensial. Karena itu, kaidah Fermat pada peminimum $\hat x$ memberikan
    \par\smallskip
    {\centering $\displaystyle 0 \in \partial f(\hat x) + K^*\partial g(K\hat x)$.\par}
    \smallskip\noindent
    Dengan kata lain, terdapat $\hat x^*\in \partial g(K\hat x)$ sedemikian sehingga $-K^*\hat x^*\in \partial f (\hat x)$. Berdasarkan~\cref{duality:lemma:subdiff_conjugate}, hal ini menyiratkan
    \begin{equation}
        \begin{aligned}
            g(K\hat x) + g^*(\hat x^*) =& \inner{\hat x^*}{K\hat x},\\
            f(\hat x) + f^*(-K^*\hat x^*)
            =& -\inner{K^*\hat x^*}{\hat x} = -\inner{\hat x^*}{K\hat x}.
        \end{aligned}
    \end{equation}
    Oleh karena itu,
    \begin{equation}
        \begin{aligned}
            \inf_x f(x) + g(Kx)
            =& f(\hat x) + g(K\hat x) \\
            =& \big(-f^*(-K^*\hat x^*) - \inner{\hat x^*}{K\hat x}\big) + \big(-g^*(\hat x^*) + \inner{\hat x^*}{K\hat x}\big)\\
            =& -f^*(-K^*\hat x^*) -g^*(\hat x^*)\\
            \leq& \sup_{x^*} -f^*(-K^*x^*) - g^*(x^*).
        \end{aligned}
    \end{equation}
    Ini membuktikan kesamaan dalam~\eqref{duality:eq:fenchel_rocka} sekaligus menunjukkan bahwa $\hat x^*$ adalah solusi masalah dual.
\end{proof}
% H07-S006 | sumber authority/habring/source-v1/duality.tex baris 248--277
% segment-id: d90.hab.v1.ch07.seg0006
\section{Optimisasi primal--dual}

\Cref{duality:eq:fenchel_duality} membuka kemungkinan beberapa metode untuk mencari solusi masalah berbentuk
\begin{equation}\label{duality:qe:problem}
    \min_x f(x) + g(Kx).
\end{equation}
Memang, berdasarkan dualitas kuat, kita dapat menangani masalah primal, masalah dual, atau masalah titik pelana
\begin{equation}\label{duality:qe:saddlepoint}
    \min_x \sup_y \Lc(x,y)\coloneqq f(x) -g^*(y) + \inner{Kx}{y}.
\end{equation}
Tujuan utama peralihan di antara formulasi primal dan dual adalah untuk \emph{memindahkan} operator $K$. Misalnya, jika $g$ tidak terdiferensialkan sehingga kita ingin menerapkan langkah proksimal pada $g$, operator $K$ dalam~\eqref{duality:qe:problem} pada umumnya membuat $\prox_{g\circ K}$ tidak mempunyai bentuk eksplisit. Sebaliknya, dalam formulasi titik pelana~\eqref{duality:qe:saddlepoint}, operator $K$ hanya muncul dalam hasil kali skalar. Selain itu, perhatikan bahwa menurut~\cref{duality:thm:Moreau:id}, proksimal suatu fungsi konjugat sama mudah atau sama sulit dihitung dengan proksimal fungsi asalnya. Lebih tepatnya, untuk $\sigma>0$ berlaku $\prox_{\sigma g^*}(y)=y-\sigma\prox_{g/\sigma}(y/\sigma)$.

Kesenjangan primal--dual didefinisikan sebagai
\begin{equation}
    \Gc(x,y) = (f(x) + g(Kx)) - (-f^*(-K^*y) - g^*(y)).
\end{equation}

\begin{lemma}
    Di bawah hipotesis dualitas kuat dan ketercapaian primal--dual yang baru saja dinyatakan, kita mempunyai $\Gc(x,y)\geq 0$, dan $\Gc(\hat x,\hat y) = 0$ jika dan hanya jika $\hat x$ menyelesaikan masalah primal dan $\hat y$ menyelesaikan masalah dual. Kesenjangan tersebut mempunyai representasi
    \begin{equation}
        \Gc(x,y) = \sup_{y'}\Lc(x,y') - \inf_{x'}\Lc(x',y).
    \end{equation}
    Khususnya, $(\hat x,\hat y)$ masing-masing optimal untuk masalah primal dan dual jika dan hanya jika, untuk setiap $x,y$,
    \begin{equation}
        \Lc(\hat x,y)\leq \Lc(\hat x,\hat y)\leq \Lc(x,\hat y).
    \end{equation}
\end{lemma}
\begin{proof}
    Latihan.
\end{proof}
% H07-S007 | sumber authority/habring/source-v1/duality.tex baris 278--335
% segment-id: d90.hab.v1.ch07.seg0007
Dengan melakukan optimisasi bergantian dalam~\eqref{duality:qe:saddlepoint}, memakai langkah proksimal terhadap $f$ dan $g^*$ serta langkah gradien terhadap hasil kali dalam, kita memperoleh aturan pembaruan berikut yang dikenal sebagai metode Arrow--Hurwicz:
\begin{equation}\label{duality:eq:AH}
    \begin{cases}
        x_{k+1} =& \prox_{\tau f}(x_k-\tau K^*y_k)\\
        y_{k+1} =& \prox_{\sigma g^*}(y_k+\sigma Kx_{k+1}).
    \end{cases}
\end{equation}
Ternyata, dalam bentuk ini metode tersebut mempunyai sifat konvergensi yang lebih buruk. Untuk memperoleh pembaruan yang lebih baik, algoritme ditulis ulang secara cermat. Dengan memakai $\prox_{\tau f}(x) = (\Id + \tau \partial f)^{-1}(x)$, dapat diperiksa secara langsung bahwa~\eqref{duality:eq:AH} ekuivalen dengan
\begin{equation}
    \begin{cases}
        x_{k+1} +\tau \partial f(x_{k+1}) =& x_k - \tau K^*y_k\\
        y_{k+1} +\sigma \partial g^*(y_{k+1}) =& y_k + \sigma Kx_{k+1}.
    \end{cases}
\end{equation}
Setelah membagi masing-masing persamaan dengan $\tau$ dan $\sigma$, serta menggunakan notasi $z=(x,y)$, bentuk ini dapat ditulis sebagai
\begin{equation}
    \begin{bmatrix}
        \frac{1}{\tau}\Id+ \partial f&0\\
        -K&\frac{1}{\sigma}\Id + \partial g^*
    \end{bmatrix}
    z_{k+1} =
    \begin{bmatrix}
        \frac{1}{\tau}\Id& - K^*\\
        0& \frac{1}{\sigma}\Id
    \end{bmatrix}
    z_k.
\end{equation}
Bahkan tanpa menelaah rincian bukti, wajar diharapkan bahwa pembaruan yang lebih \emph{simetris} akan lebih baik, khususnya karena sifat menguntungkan matriks simetris. Karena itu, kita mempertimbangkan skema berikut, yang tidak mengubah titik tetap yang mungkin:
\begin{equation}
    \begin{bmatrix}
        \frac{1}{\tau}\Id+ \partial f&0\\
        -2K&\frac{1}{\sigma}\Id + \partial g^*
    \end{bmatrix}
    z_{k+1} =
    \begin{bmatrix}
        \frac{1}{\tau}\Id& -K^*\\
        -K & \frac{1}{\sigma}\Id
    \end{bmatrix}
    z_k.
\end{equation}
Matriks pada ruas kanan kini simetris dan definit positif selama $\sigma\tau \|K\|^2<1$.
Kembali ke operator proksimal, metode gradien hibrida primal--dual ini, yang sering disebut metode Chambolle--Pock berdasarkan~\cite{chambolle2011first}, berbentuk
\begin{equation}\label{duality:eq:pdhg}
    \begin{cases}
        x_{k+1} =& \prox_{\tau f}(x_k-\tau K^*y_k)\\
        y_{k+1} =& \prox_{\sigma g^*}(y_k+\sigma K(2x_{k+1}-x_k)).
    \end{cases}
\end{equation}

Untuk analisis, PDHG akan ditulis dalam bentuk abstrak
\begin{equation}\label{duality:eq:pdhg2}
    \begin{cases}
        x^+ =& \arg\min_x \frac{1}{2\tau}|x-x^-|^2 + \inner{x}{K^*\bar y} + f(x)\\
        y^+ =& \arg\min_y \frac{1}{2\sigma}|y-y^-|^2 - \inner{y}{K\bar x} + g^*(y).
    \end{cases}
\end{equation}
Di sini, superskrip plus menyatakan variabel yang telah diperbarui, superskrip minus menyatakan iterasi lama, sedangkan garis atas masing-masing menyatakan nilai $x$ yang digunakan untuk pembaruan $y$ dan nilai $y$ yang digunakan untuk pembaruan $x$.
% H07-S008 | sumber authority/habring/source-v1/duality.tex baris 336--394
% segment-id: d90.hab.v1.ch07.seg0008
\begin{lemma}
    Misalkan $X$ dan $Y$ ruang Hilbert riil berdimensi hingga, $f\in\Gamma_0(X)$, $g\in\Gamma_0(Y)$, $K:X\rightarrow Y$ linear terbatas, serta $\tau,\sigma>0$. Aturan pembaruan~\eqref{duality:eq:pdhg2} memenuhi
    \begin{equation}
        \begin{aligned}
            \Lc(x^+,y) - \Lc(x,y^+)\leq& \frac{1}{2\tau}|x-x^-|^2 + \frac{1}{2\sigma}|y-y^-|^2\\
            &-\frac{1}{2\tau}|x-x^+|^2 - \frac{1}{2\sigma}|y-y^+|^2
            -\frac{1}{2\tau}|x^+-x^-|^2 - \frac{1}{2\sigma}|y^+-y^-|^2\\
            &+\inner{x^+-x}{K^*(y-\bar{y})}  - \inner{y^+-y}{K(x-\bar{x})}.
        \end{aligned}
    \end{equation}
\end{lemma}
\begin{proof}
    Berdasarkan kekonveksan kuat, pembaruan $x^+$ dan $y^+$ memenuhi
    \begin{equation}
        \begin{aligned}
            \frac{1}{2\tau}|x^+-x^-|^2 + \inner{x^+}{K^*\bar y} + f(x^+) + \frac{1}{2\tau}|x-x^+|^2
            \leq& \frac{1}{2\tau}|x-x^-|^2 + \inner{x}{K^*\bar y} + f(x)\\
            \frac{1}{2\sigma}|y^+-y^-|^2 - \inner{y^+}{K\bar x} + g^*(y^+) + \frac{1}{2\sigma}|y-y^+|^2
            \leq& \frac{1}{2\sigma}|y-y^-|^2 - \inner{y}{K\bar x} + g^*(y).
        \end{aligned}
    \end{equation}
    Menjumlahkan kedua ketaksamaan dan menyusun ulang suku-sukunya memberikan hasil yang diinginkan.
\end{proof}

\begin{theorem}[Konvergensi PDHG]
    Misalkan $X$ dan $Y$ ruang Hilbert riil berdimensi hingga, $f\in\Gamma_0(X)$, $g\in\Gamma_0(Y)$, $K:X\rightarrow Y$ linear terbatas, serta $\tau,\sigma>0$. Perhatikan pembaruan~\eqref{duality:eq:pdhg} dan andaikan ukuran langkah memenuhi $\tau\sigma\leq \|K\|^{-2}$; untuk $K=0$, syarat ini dipahami otomatis terpenuhi. Untuk $N\geq1$, definisikan rerata ergodik $X_N = \frac{1}{N}\sum_{k=1}^{N}x_k$ dan $Y_N = \frac{1}{N}\sum_{k=1}^{N}y_k$. Maka
    \begin{equation}
        \Lc(X_{N},y) - \Lc(x,Y_{N})\leq \frac{1}{N}\bigg[\frac{1}{2\tau}|x-x_0|^2 + \frac{1}{2\sigma}|y-y_0|^2 - \inner{y-y_0}{K(x-x_{0})}\bigg].
    \end{equation}
\end{theorem}
\begin{rem}
    Perhatikan bahwa ukuran langkah sama sekali tidak bergantung pada fungsi $f$ dan $g$, melainkan hanya pada operator $K$.

\end{rem}
\begin{proof}
    {\small
    \begin{equation}
        \begin{aligned}
            \Lc(x_{k+1},y) - \Lc(x,y_{k+1})
            \leq& \frac{1}{2\tau}|x-x_k|^2 + \frac{1}{2\sigma}|y-y_k|^2\\
            &-\frac{1}{2\tau}|x-x_{k+1}|^2 - \frac{1}{2\sigma}|y-y_{k+1}|^2
            -\frac{1}{2\tau}|x_{k+1}-x_k|^2 - \frac{1}{2\sigma}|y_{k+1}-y_k|^2\\
            &+\inner{x_{k+1}-x}{K^*(y-y_k)}  - \inner{y_{k+1}-y}{K(x-x_{k+1})}\\
            &\quad + \inner{y_{k+1}-y}{K(x_{k+1}-x_k)}\\
            \leq& \bigg[\frac{1}{2\tau}|x-x_k|^2 + \frac{1}{2\sigma}|y-y_k|^2 - \inner{y-y_k}{K(x-x_{k})}\bigg]\\
            &-\bigg[\frac{1}{2\tau}|x-x_{k+1}|^2 + \frac{1}{2\sigma}|y-y_{k+1}|^2 - \inner{y-y_{k+1}}{K(x-x_{k+1})}\bigg]\\
            &-\bigg[ \frac{1}{2\tau}|x_{k+1}-x_k|^2 + \frac{1}{2\sigma}|y_{k+1}-y_k|^2\\
            &\qquad - \inner{x_{k+1}-x_k}{K^*(y_{k+1}-y_{k})}\bigg].
        \end{aligned}
    \end{equation}
    }
    Jika $K=0$, suku campuran dalam setiap tanda kurung siku lenyap, sehingga semua ekspresi itu tidak negatif. Jika $K\neq0$, nonnegativitas mengikuti
    \par\smallskip
    {\centering\small
    $\displaystyle \inner{Kx}{y}\leq \|Kx\|\|y\|\leq\|K\| \big(\frac{\delta}{2}\|x\|^2 + \frac{1}{2\delta}\|y\|^2\big)\leq \frac{1}{2\tau}\|x\|^2 + \frac{1}{2\sigma}\|y\|^2$,
    \par}
    \smallskip
    dengan ketaksamaan terakhir diperoleh dari pilihan $\delta = \frac{1}{\tau\|K\|}$ dan fakta bahwa $\|K\|^2\leq \frac{1}{\tau\sigma}$.
    Penjumlahan terhadap $k$ menghasilkan
    \begin{equation}
        \sum_{k=0}^{N-1} \Lc(x_{k+1},y) - \Lc(x,y_{k+1}) \leq \bigg[\frac{1}{2\tau}|x-x_0|^2 + \frac{1}{2\sigma}|y-y_0|^2 - \inner{y-y_0}{K(x-x_{0})}\bigg].
    \end{equation}
    Karena pemetaan $(x,y)\mapsto \Lc(x,y') - \Lc(x',y)$ konveks untuk $x',y'$ tetap, kita memperoleh
    \begin{equation}
        \resizebox{0.88\linewidth}{!}{$\displaystyle
        \Lc(X_{N},y) - \Lc(x,Y_{N})\leq \frac{1}{N}\sum_{k=0}^{N-1} \big(\Lc(x_{k+1},y) - \Lc(x,y_{k+1})\big) \leq \frac{1}{N}\bigg[\frac{1}{2\tau}|x-x_0|^2 + \frac{1}{2\sigma}|y-y_0|^2 - \inner{y-y_0}{K(x-x_{0})}\bigg].
        $}
    \end{equation}
\end{proof}

% H07-S009 | sumber authority/habring/source-v1/duality.tex baris 395--432
% segment-id: d90.hab.v1.ch07.seg0009
\section{Metode pengali arah bergantian (ADMM)}
Misalkan $f,g\in\Gamma_0(\R^d)$, $A,B:\R^d\rightarrow\R^m$ linear, dan $z\in\R^m$. Perhatikan masalah optimisasi berkendala
\begin{equation}\label{eq:admm1}
    \min_{Ax + By = z} f(x) + g(y).
\end{equation}
Masalah ini dapat dirumuskan ulang sebagai masalah titik pelana tak berkendala untuk \emph{Lagrangian teraugmentasi}; yaitu, untuk setiap $\gamma>0$, kita mempertimbangkan
\begin{equation}\label{eq:admm2}
    \min_{x,y}\sup_{\lambda} L_\gamma (x,y,\lambda) \coloneqq f(x) + g(y) + \inner{\lambda}{Ax + By - z} + \frac{\gamma}{2}\|Ax + By - z\|^2.
\end{equation}
Pernyataan berikut dapat ditunjukkan secara langsung.
\begin{lemma}
    Masalah~\eqref{eq:admm1} dan~\eqref{eq:admm2} ekuivalen.
\end{lemma}
\begin{proof}
    Latihan!
\end{proof}
Salah satu cara langsung untuk menyelesaikan masalah titik pelana tersebut adalah mengoptimalkan secara bergantian terhadap $x$ dan $y$, lalu melakukan langkah kenaikan gradien terhadap $\lambda$:
\begin{equation}\label{duality:eq:admm_algo}
    \begin{cases}
        x_{k+1} = \arg\min_{x} f(x) + \inner{\lambda_k}{Ax} + \frac{\gamma}{2}\|Ax + By_k - z\|^2\\
        y_{k+1} = \arg\min_{y} g(y) + \inner{\lambda_k}{By} + \frac{\gamma}{2}\|Ax_{k+1} + By - z\|^2\\
        \lambda_{k+1} = \lambda_k + \gamma (Ax_{k+1} + By_{k+1} - z).
    \end{cases}
\end{equation}
Perhatikan bahwa $\gamma$ digunakan sebagai ukuran langkah untuk kenaikan gradien. Pilihan ini dapat dimotivasi secara intuitif oleh fakta bahwa, berdasarkan definisi pembaruan,
\begin{equation}\label{duality:eq:admm3}
    \begin{aligned}
        0 \in \partial g(y_{k+1}) + B^* \lambda_k + \gamma B^*(Ax_{k+1} + By_{k+1}-z).
    \end{aligned}
\end{equation}
Dengan mengalikan pembaruan $\lambda$ oleh $B^*$ dan memasukkannya ke~\eqref{duality:eq:admm3}, diperoleh
\begin{equation}
    B^*\lambda_{k+1} \in -\partial g(y_{k+1}).
\end{equation}
Dengan demikian, tripel ADMM yang benar-benar diperbarui melalui~\eqref{duality:eq:admm_algo} memenuhi $0\in\partial_y L_0(x_{k+1},y_{k+1},\lambda_{k+1})$. Pembaruan bergantian tersebut pada umumnya tidak memberikan stasioneritas terhadap $x$, karena submasalah $x$ masih memakai $y_k$. Hanya dalam skema hipotetis yang meminimumkan $x$ dan $y$ secara bersama kondisi serupa diperoleh terhadap $x$, sehingga tripel hasil pembaruan bersama itu memenuhi stasioneritas primal penuh untuk $L_0$.
% H07-S010 | sumber authority/habring/source-v1/duality.tex baris 433--502
% segment-id: d90.hab.v1.ch07.seg0010
Kita dapat menunjukkan hasil konvergensi dasar berikut untuk ADMM.
\begin{theorem}
    Ambil $\gamma>0$. Andaikan terdapat suatu titik pelana dari $L_0(x,y,\lambda)$ dan, pada setiap iterasi, kedua submasalah minimisasi dalam algoritme ADMM di atas mencapai minimumnya; pilih sembarang peminimum tersebut. Dengan $r_k\coloneq Ax_k+By_k-z$, berlaku $r_k\rightarrow0$, $B(y_{k+1}-y_k)\rightarrow0$, dan
    \begin{equation}
        f(x_k)+g(y_k) \rightarrow \min_{Ax + By = z} f(x) + g(y).
    \end{equation}
\end{theorem}
\begin{proof}
    Bukti menggunakan suatu teknik penting, yaitu \emph{fungsional Lyapunov}. Kita mendefinisikan fungsional $V(y,\lambda)$ yang tidak negatif (terbatas dari bawah) dan menurun sepanjang iterasi. Pada dasarnya, fungsional Lyapunov menggantikan fungsi objektif ketika penurunan objektif itu sendiri tidak dapat dibuktikan. Dalam kasus ini, kita menggunakan
    \begin{equation}
        V(y,\lambda) = \gamma \|B(y-y^*)\|^2 + \gamma^{-1} \|\lambda-\lambda^*\|^2,
    \end{equation}
    dengan $(x^*,y^*,\lambda^*)$ suatu titik pelana bagi $L_0$. Penurunan fungsional Lyapunov akan memberikan konvergensi melalui argumen penjumlahan teleskopik yang lazim. Akan tetapi, pembuktian penurunan ini memerlukan beberapa taksiran teknis. Demi kesederhanaan, tuliskan residu $r_k \coloneq Ax_k + By_k-z$ dan nilai objektif $p_k \coloneq f(x_k) + g(y_k)$; gunakan pula $r^*$ dan $p^*$ untuk titik pelana. Bukti dibagi menjadi tiga langkah.
    \paragraph{Langkah 1}
    Karena $(x^*,y^*,\lambda^*)$ merupakan titik pelana bagi $L_0$, untuk setiap $k$ berlaku $L_0(x^*,y^*,\lambda^*)\leq L_0(x_{k+1},y_{k+1},\lambda^*)$. Pernyataan ini ekuivalen dengan
    \begin{equation}\label{duality:eq:proof_admm1}
        p^*-p_{k+1}\leq \inner{\lambda^*}{r_{k+1}}.
    \end{equation}
    \paragraph{Langkah 2}
    Di sisi lain, berdasarkan optimalitas $x_{k+1}$ dan $y_{k+1}$ dalam~\eqref{duality:eq:admm_algo}, kita memperoleh
    \begin{equation}\label{duality:eq:proof_admm2}
        \begin{aligned}
            0 &\in \partial f(x_{k+1}) + A^*\lambda_k + \gamma A^*(r_{k+1} + B(y_k-y_{k+1}))\\
            0 &\in \partial g(y_{k+1}) + B^*\lambda_k + \gamma B^*r_{k+1}.
        \end{aligned}
    \end{equation}
    Selain itu, pembaruan $\lambda$ dapat ditulis sebagai $\lambda_{k+1} = \lambda_k + \gamma r_{k+1}$, sehingga
    \begin{equation}
        \begin{aligned}
            A^* \lambda_{k+1} &= A^*\lambda_k + \gamma A^*r_{k+1}\\
            B^* \lambda_{k+1} &= B^*\lambda_k + \gamma B^*r_{k+1}.
        \end{aligned}
    \end{equation}
    Ketika dimasukkan ke~\eqref{duality:eq:proof_admm2}, persamaan ini memberikan
    \begin{equation}
        \begin{aligned}
            0 &\in \partial f(x_{k+1}) + A^*(\lambda_{k+1} + \gamma B(y_k-y_{k+1}))\\
            0 &\in \partial g(y_{k+1}) + B^*\lambda_{k+1}.
        \end{aligned}
    \end{equation}
    Hal ini menyiratkan bahwa $x_{k+1}$ meminimumkan
    \begin{equation}
        x\mapsto f(x) + \inner{A x}{\lambda_{k+1} + \gamma B(y_k-y_{k+1})},
    \end{equation}
    sedangkan $y_{k+1}$ meminimumkan
    \begin{equation}
        y\mapsto g(y) + \inner{By}{\lambda_{k+1}}.
    \end{equation}
    Akibatnya, dari optimalitas diperoleh
    \begin{equation}
        \begin{aligned}
            f(x_{k+1}) + &\inner{A x_{k+1}}{\lambda_{k+1} + \gamma B(y_k-y_{k+1})} + g(y_{k+1}) + \inner{By_{k+1}}{\lambda_{k+1}}\\
            &\leq
            f(x^*) + \inner{Ax^*}{\lambda_{k+1} + \gamma B(y_k-y_{k+1})} + g(y^*) + \inner{By^*}{\lambda_{k+1}}.
        \end{aligned}
    \end{equation}
    Karena titik pelana pasti memenuhi $Ax^* + By^* = z$, sedikit penyusunan ulang memberikan
    \begin{equation}
        \begin{aligned}
            p_{k+1} - p^*
            \leq \gamma \inner{A (x^*-x_{k+1})}{B(y_k-y_{k+1})} - \inner{r_{k+1}}{\lambda_{k+1}}.
        \end{aligned}
    \end{equation}
    Terakhir, kita dapat memasukkan $A(x^* - x_{k+1}) = z - By^* - (r_{k+1} + z - By_{k+1}) = -r_{k+1} + B(y_{k+1}-y^*)$ untuk memperoleh
    \begin{equation}\label{duality:eq:proof_admm3}
        \begin{aligned}
            p_{k+1} - p^*
            \leq -\gamma \inner{-r_{k+1} + B(y_{k+1}-y^*)}{B(y_{k+1}-y_k)} - \inner{r_{k+1}}{\lambda_{k+1}}.
        \end{aligned}
    \end{equation}
% H07-S011 | sumber authority/habring/source-v1/duality.tex baris 503--597
% segment-id: d90.hab.v1.ch07.seg0011
    \paragraph{Langkah 3}
    Dengan menjumlahkan~\eqref{duality:eq:proof_admm1} dan~\eqref{duality:eq:proof_admm3}, lalu mengalikan dengan $2$, diperoleh
    \begin{equation}\label{duality:eq:proof_admm4}
        -2\gamma \inner{r_{k+1}}{B(y_{k+1}-y_k)} + 2\gamma\inner{B(y_{k+1}-y^*)}{B(y_{k+1}-y_k)} + 2\inner{r_{k+1}}{\lambda_{k+1}-\lambda^*}\leq 0.
    \end{equation}
    Kita akan mengolah ketaksamaan ini. Suku ketiga dapat ditulis ulang menggunakan $\lambda_{k+1} = \lambda_k + \gamma r_{k+1}$ sebagai
    \begin{equation}
        2\inner{r_{k+1}}{\lambda_{k+1}-\lambda^*}
        = \gamma\|r_{k+1}\|^2 + \gamma\|r_{k+1}\|^2 + 2\inner{r_{k+1}}{\lambda_k -\lambda^*}.
    \end{equation}
    Dengan mengganti $r_{k+1} = (\lambda_{k+1}-\lambda_k)/\gamma$ pada dua suku terakhir, diperoleh
    \begin{equation}
        \gamma\|r_{k+1}\|^2 + \gamma^{-1}\|\lambda_{k+1}-\lambda_k\|^2 + 2\gamma^{-1}\inner{\lambda_{k+1}-\lambda_k}{\lambda_k -\lambda^*}.
    \end{equation}
    Dengan mengganti $\lambda_{k+1}-\lambda_k = \lambda_{k+1}-\lambda^* + \lambda^* - \lambda_k$ di atas, bentuk ini ekuivalen dengan
    \begin{equation}
        \gamma\|r_{k+1}\|^2 + \gamma^{-1}\big(\|\lambda_{k+1}-\lambda^*\|^2 - \|\lambda_{k}-\lambda^*\|^2\big).
    \end{equation}
    Jadi,~\eqref{duality:eq:proof_admm4} ekuivalen dengan
    \begin{equation}\label{duality:eq:proof_admm5}
        \begin{aligned}
            0\geq&-2\gamma \inner{r_{k+1}}{B(y_{k+1}-y_k)} + 2\gamma\inner{B(y_{k+1}-y^*)}{B(y_{k+1}-y_k)} + \gamma\|r_{k+1}\|^2 \\
            &+ \gamma^{-1}\big(\|\lambda_{k+1}-\lambda^*\|^2 - \|\lambda_{k}-\lambda^*\|^2\big).
        \end{aligned}
    \end{equation}
    Sekarang kita tulis ulang tiga suku pertama,
    \begin{equation}\label{duality:eq:proof_admm6}
        -2\gamma \inner{r_{k+1}}{B(y_{k+1}-y_k)} + 2\gamma\inner{B(y_{k+1}-y^*)}{B(y_{k+1}-y_k)} + \gamma\|r_{k+1}\|^2.
    \end{equation}
    Dengan memperhatikan bahwa
    \begin{equation}
        -2\inner{r_{k+1}}{B(y_{k+1}-y_k)} + \|r_{k+1}\|^2
        = \|r_{k+1} - B(y_{k+1}-y_k)\|^2 - \|B(y_{k+1}-y_k)\|^2,
    \end{equation}
    kita mendapatkan bahwa~\eqref{duality:eq:proof_admm5} ekuivalen dengan
    \begin{equation}
        \gamma \|r_{k+1} - B(y_{k+1}-y_k)\|^2 - \gamma\|B(y_{k+1}-y_k)\|^2 + 2\gamma\inner{B(y_{k+1}-y^*)}{B(y_{k+1}-y_k)},
    \end{equation}
    yang ekuivalen dengan
    \begin{equation}
        \gamma \|r_{k+1} - B(y_{k+1}-y_k)\|^2 + \gamma\big( \|B(y_{k+1}-y^*)\|^2 - \|B(y_{k}-y^*)\|^2\big).
    \end{equation}
    Secara keseluruhan,~\eqref{duality:eq:proof_admm5} kemudian menghasilkan
    {\small
    \begin{equation}
        0\geq \gamma \|r_{k+1} - B(y_{k+1}-y_k)\|^2 + \gamma\big( \|B(y_{k+1}-y^*)\|^2 - \|B(y_{k}-y^*)\|^2\big) + \gamma^{-1}\big(\|\lambda_{k+1}-\lambda^*\|^2 - \|\lambda_{k}-\lambda^*\|^2\big).
    \end{equation}
    }
    Terakhir, perhatikan bahwa
    \begin{equation}
        \|r_{k+1} - B(y_{k+1}-y_k)\|^2
        =\|r_{k+1}\|^2 + \|B(y_{k+1}-y_k)\|^2 - 2\inner{r_{k+1}}{B(y_{k+1}-y_k)}.
    \end{equation}
    Kita ingin menunjukkan bahwa suku
    \begin{equation}
        \inner{r_{k+1}}{B(y_{k+1}-y_k)} \leq 0.
    \end{equation}
    Untuk $k\geq 1$, ingat bahwa $y_k$ meminimumkan $g(y) + \inner{By}{\lambda_k}$. Dengan demikian,
    \begin{equation}
        \begin{aligned}
            g(y_k) + \inner{By_k}{\lambda_k}
            &\leq g(y_{k+1}) + \inner{By_{k+1}}{\lambda_k}\\
            g(y_{k+1}) + \inner{By_{k+1}}{\lambda_{k+1}}
            &\leq g(y_k) + \inner{By_k}{\lambda_{k+1}}.
        \end{aligned}
    \end{equation}
    Penjumlahan kedua ketaksamaan memberikan
    \begin{equation}
        \begin{aligned}
            \inner{B(y_{k+1}-y_k)}{\lambda_{k+1}-\lambda_k}
            \leq 0.
        \end{aligned}
    \end{equation}
    Dengan memasukkan $\lambda_{k+1}-\lambda_k = \gamma r_{k+1}$, diperoleh tanda yang diinginkan. Oleh karena itu,
    \begin{equation}
        V(y_{k+1},\lambda_{k+1}) + \gamma\|r_{k+1}\|^2 + \gamma\|B(y_{k+1}-y_k)\|^2 \leq V(y_{k},\lambda_{k}),\qquad k\geq1.
    \end{equation}
    Menjumlahkan terhadap $k$ menghasilkan, untuk $K\geq2$,
    \begin{equation}
        \sum_{k=1}^{K-1}\gamma\big(\|r_{k+1}\|^2 + \|B(y_{k+1}-y_k)\|^2\big) \leq V(y_{1},\lambda_{1}).
    \end{equation}
    Dari penurunan fungsional Lyapunov, segera diperoleh bahwa $\lambda_k$ dan $By_k$ merupakan barisan terbatas.
    Selain itu,
    \begin{equation}\label{duality:eq:proof_admm7}
        \begin{cases}
            r_k&\rightarrow 0\\
            B(y_{k+1}-y_k)&\rightarrow 0.
        \end{cases}
    \end{equation}
    Berdasarkan~\eqref{duality:eq:proof_admm1} dan~\eqref{duality:eq:proof_admm3}, kita mempunyai
    \begin{equation}
        \begin{aligned}
            -\inner{\lambda^*}{r_{k+1}}\leq p_{k+1}-p^*\leq -\gamma\inner{-r_{k+1} + B(y_{k+1}-y^*)}{B(y_{k+1}-y_k)} - \inner{r_{k+1}}{\lambda_{k+1}}.
        \end{aligned}
    \end{equation}
    Kemudian~\eqref{duality:eq:proof_admm7} dan keterbatasan $By_k$ menyiratkan bahwa $p_{k}\rightarrow p^*$.
\end{proof}


\backmatter
\printbibliography[heading=bibintoc,title={Daftar pustaka}]

\chapter*{Catatan koreksi terhadap sumber}
Terjemahan ini mempertahankan isi sumber kecuali pada koreksi yang dapat ditentukan berikut. Koreksi dilakukan oleh penyusun edisi bahasa Indonesia dan bukan perubahan yang disahkan oleh penulis sumber. Rekaman terstruktur dengan identitas peristiwa O015-HAB-ADV-0050 sampai O015-HAB-ADV-0075 menyertai sumber edisi ini.

\begin{enumerate}
    \item Pada contoh konjugat, urutan pasangan dual dibetulkan menjadi $\inner{x^*}{x}$, dan konjugat indikator bola satuan dual diidentifikasi sebagai norma asal pada $V$ setelah identifikasi $V^{**}\cong V$.
    \item Lema keproperan konjugat kini menyatakan hipotesis keberadaan sedikitnya satu subgradien yang dipakai buktinya; evaluasi yang salah tipe $f^*(\hat x)$ dibetulkan menjadi $f^*(g)$, dan fungsi dalam supremum disebut afin kontinu.
    \item Bukti bikonjugat memakai $f(x)$ ketika memilih $y=x$ dan memulihkan dua tanda ketaksamaan yang hilang pada cabang $\alpha=0$.
    \item Identitas Moreau dinyatakan di bawah konvensi ruang Hilbert bab ini agar pemetaan proksimal yang dipakai memang terdefinisi dan bernilai tunggal dalam lingkup bersama yang dinyatakan.
    \item Rumus kesenjangan primal--dual diperbaiki agar konsisten tipe dan tanda dengan dual Fenchel: $f(x)+g(Kx)-[-f^*(-K^*y)-g^*(y)]$.
    \item Pembaruan dual PDHG memakai $\prox_{\sigma g^*}$; submasalah abstraknya diminimumkan terhadap $y$; dan koefisien sisa kekonveksan kuat untuk $y$ memakai $1/(2\sigma)$.
    \item Rerata ergodik PDHG dan langkah penjumlahannya memakai iterasi $1$ sampai $N$, sesuai ketaksamaan satu langkah untuk $x_{k+1}$ dan $y_{k+1}$.
    \item Submasalah kedua ADMM diminimumkan terhadap $y$. Notasi gradien yang tidak berlaku bagi $g$ nonsmooth diganti dengan subdiferensial. Tripel ADMM aktual hanya dinyatakan stasioner terhadap $y$; stasioneritas primal penuh dibatasi pada skema hipotetis dengan minimisasi bersama.
    \item Substitusi residu ADMM dibetulkan menjadi $Ax_{k+1}=r_{k+1}+z-By_{k+1}$ sebelum menurunkan $A(x^*-x_{k+1})$.
    \item Dalam penurunan Lyapunov ADMM, tanda suku silang, tanda yang dibuktikan oleh optimalitas, faktor $\gamma$, dan indeks awal $k\geq1$ dibetulkan. Penjumlahan teleskopik dimulai dari iterasi pertama yang memenuhi kedua kondisi optimalitas $y$.
    \item Batas atas terakhir untuk galat nilai objektif mengembalikan tanda minus dan faktor $\gamma$ yang sudah terdapat pada ketaksamaan berlabel sebelumnya.
    \item Kemunculan kedua label \texttt{duality:eq:proof\_admm6} dipetakan secara transparan ke \texttt{duality:eq:proof\_admm7}; rujukan tunggalnya diperbarui bersama-sama.
    \item Salah eja, frasa ganda, dan karakter kendali berlebih dalam prosa sumber dinormalkan dalam terjemahan tanpa mengubah matematika lain.
    \item Lingkup bab dibuat koheren sebagai ruang Hilbert riil berdimensi hingga dengan dual kontinu dan identifikasi Riesz; $X$ dan $Y$ dinyatakan eksplisit. Kelas $\Gamma_0$ memakai konvensi fungsi proper, konveks, semikontinu bawah dengan nilai dalam $(-\infty,+\infty]$.
    \item Karakterisasi regularisasi $\Gamma$ dibatasi pada keluarga minoran proper yang tidak kosong; jika keluarga itu kosong, supremumnya secara eksplisit adalah fungsi yang identik $-\infty$.
    \item Dalam bukti Fenchel--Rockafellar, inklusi optimalitas kini dijustifikasi oleh kaidah Fermat serta aturan jumlah dan rantai di bawah syarat kualifikasi interior yang telah dinyatakan.
    \item Maksimum titik demi titik yang ketercapaiannya tidak dijamin diganti dengan supremum, dan karakterisasi kesenjangan nol dinyatakan di bawah hipotesis dualitas kuat serta ketercapaian primal--dual.
    \item Identitas Moreau berskala $\prox_{\sigma g^*}(y)=y-\sigma\prox_{g/\sigma}(y/\sigma)$ ditambahkan untuk menjelaskan langkah proksimal dual PDHG.
    \item Teorema PDHG kini menyatakan fungsi dalam $\Gamma_0$, ruang dan operatornya, serta positifnya ukuran langkah; penghitung iterasi diganti dari $K$ menjadi $N$ agar tidak bertabrakan dengan operator $K$, dan kasus operator nol dinyatakan otomatis memenuhi syarat ukuran langkah.
    \item Bagian ADMM kini memberi tipe $A,B:\R^d\to\R^m$ dan $z\in\R^m$, mengasumsikan $\gamma>0$ serta ketercapaian setiap submasalah eksak, dan menyatakan hanya konvergensi residu, perubahan citra blok, serta nilai objektif yang benar-benar dibuktikan.
    \item Himpunan $C$ pada contoh fungsi indikator disyaratkan tak kosong agar $\delta_C$ proper dan fungsi pendukungnya tetap berada dalam kodomain nilai diperluas yang dinyatakan.
    \item Tanda $\alpha\geq0$ dalam bukti Fenchel--Moreau kini diperoleh dari suatu $\bar x\in\dom(f)$ yang dijamin oleh keproperan, lalu dengan membiarkan tinggi epigraf menuju tak hingga; titik terpisah $x_0$ tidak lagi diasumsikan berada dalam domain efektif.
    \item Bukti PDHG menangani operator nol secara terpisah, ketika suku campuran lenyap, dan hanya memakai $\delta=1/(\tau\|K\|)$ ketika $K\neq0$.
    \item Lema satu langkah PDHG kini menyatakan sendiri hipotesis ruang Hilbert berdimensi hingga, fungsi dalam $\Gamma_0$, operator linear terbatas, dan ukuran langkah positif yang diperlukan karena lema itu mendahului teorema konvergensi.
    \item Bikonjugasi dan ketaksamaan $f^{**}\leq f$ kini mensyaratkan $f$ serta $f^*$ proper. Dalam kasus keluarga minoran $\Gamma$ kosong, hanya supremum ruas kanan yang didefinisikan sebagai fungsi identik $-\infty$; di bawah konvensi masukan proper bab ini, fungsi tersebut tidak dinotasikan sebagai $f^{**}$.
\end{enumerate}

\chapter*{Atribusi, lisensi, dan nondukungan}
Karya sumber: Andreas Habring, \emph{Lecture Notes: Convex Optimization}, arXiv:2607.11664v1, 13 Juli 2026, \url{https://arxiv.org/abs/2607.11664v1}. Karya sumber tersedia berdasarkan Creative Commons Attribution 4.0 International (CC BY 4.0), \url{https://creativecommons.org/licenses/by/4.0/}.

Edisi bahasa Indonesia ini merupakan karya turunan. Hak cipta materi sumber tetap pada pemegang haknya. Anda boleh membagikan dan mengadaptasinya dengan atribusi yang sesuai, tautan lisensi, dan penandaan perubahan sebagaimana disyaratkan CC BY 4.0.

Andreas Habring dan TU Graz tidak menyusun, memeriksa, menyetujui, atau mendukung terjemahan ini. Penyebutan nama mereka semata-mata untuk atribusi karya sumber.

\end{document}
